\chapter{The Cell Cycle and Mitosis}\label{ch:g11:cell-cycle-mitosis}

A growing tissue is not a bag of identical resting \omterm{def:g10:the-cell:cell}{cells}: \omterm{def:g10:the-cell:cell}{cells} grow, copy their
\omterm{def:g11:dna-replication:helix}{DNA}, and divide. The ordered programme of \omterm{def:g10:scientific-biology:properties}{growth} and division is the
\emph{\omterm{def:g11:cell-cycle-mitosis:cycle}{cell cycle}}; the division of the nucleus that preserves \omterm{def:g11:dna-replication:chromosome}{chromosome} number
in somatic \omterm{def:g10:the-cell:cell}{cells} is \emph{\omterm{def:g11:cell-cycle-mitosis:mitosis}{mitosis}}. This chapter names the phases, the structure
of duplicated \omterm{def:g11:dna-replication:chromosome}{chromosomes}, and the idea of \omterm{def:g11:cell-cycle-mitosis:checkpoint}{checkpoints}.

\section{The cell cycle}

\begin{definition}[Cell cycle]\label{def:g11:cell-cycle-mitosis:cycle}
The \emph{cell cycle}\index{cell cycle} is the \omterm{def:g10:molecules-of-life:proteins}{sequence} of stages a \omterm{def:g10:the-cell:cell}{cell} passes
through from one division to the next. In \omterm{def:g10:the-cell:prok-euk}{eukaryotes} it comprises
\emph{interphase}\index{interphase} (G$_1$, \omterm{def:g11:cell-cycle-mitosis:interphase}{S}, G$_2$) and the
\emph{mitotic phase}\index{mitotic phase} (\omterm{def:g11:cell-cycle-mitosis:mitosis}{mitosis} + \omterm{def:g11:cell-cycle-mitosis:cytokinesis}{cytokinesis}).
\end{definition}

\begin{definition}[Interphase stages]\label{def:g11:cell-cycle-mitosis:interphase}
\begin{itemize}
  \item \emph{G$_1$}\index{G1 phase} (first gap): the \omterm{def:g10:the-cell:cell}{cell} grows, performs its
        functions, and decides whether to commit to division.
  \item \emph{S}\index{S phase} (synthesis): \omterm{def:g11:dna-replication:helix}{DNA} is replicated
        (\cref{ch:g11:dna-replication}); each \omterm{def:g11:dna-replication:chromosome}{chromosome} becomes two \omterm{def:g11:cell-cycle-mitosis:chromatid}{sister
        chromatids}.
  \item \emph{G$_2$}\index{G2 phase} (second gap): further \omterm{def:g10:scientific-biology:properties}{growth} and preparation
        for \omterm{def:g11:cell-cycle-mitosis:mitosis}{mitosis}; many \omterm{def:g10:the-cell:prok-euk}{organelles} are duplicated.
\end{itemize}
Some \omterm{def:g10:the-cell:cell}{cells} exit the cycle into a resting state called \emph{G$_0$}\index{G0 phase}.
\end{definition}

\begin{omfigure}
\begin{tikzpicture}[font=\small]
  \draw[omDef, line width=3pt] (0,0) circle (2.0);
  % arcs labels
  \node at (90:2.55) {G$_1$};
  \node at (0:2.55) {S};
  \node at (-90:2.55) {G$_2$};
  \node at (180:2.7) {M};
  \draw[omMeth, line width=4pt]
    (120:2.0) arc[start angle=120, end angle=30, radius=2.0];
  \draw[omProp, line width=4pt]
    (30:2.0) arc[start angle=30, end angle=-60, radius=2.0];
  \draw[omThm, line width=4pt]
    (-60:2.0) arc[start angle=-60, end angle=-120, radius=2.0];
  \draw[omExo, line width=4pt]
    (-120:2.0) arc[start angle=-120, end angle=-240, radius=2.0];
  \node[align=center] at (0,0) {cell\\cycle};
  \draw[-Stealth, thick] (0.3,0.3) -- (0.9,0.9);
\end{tikzpicture}

{\small The \omterm{def:g10:the-cell:prok-euk}{eukaryotic cell} cycle: long \omterm{def:g11:cell-cycle-mitosis:cycle}{interphase} (G$_1$--S--G$_2$), short M
phase. Relative durations vary by \omterm{def:g10:the-cell:cell}{cell} type; embryonic cycles can omit gaps.}
\end{omfigure}

\begin{proposition}[Why S comes before M]\label{prop:g11:cell-cycle-mitosis:order}
\omterm{def:g11:cell-cycle-mitosis:mitosis}{Mitosis} segregates already-duplicated \omterm{def:g11:dna-replication:chromosome}{chromosomes}. If division occurred without
\omterm{def:g11:cell-cycle-mitosis:interphase}{S phase}, each daughter would receive an incomplete \omterm{prop:g11:dna-replication:genome}{genome}. The cycle \omterm{def:g10:biodiversity-classification:hierarchy}{order}
\omterm{def:g10:scientific-biology:properties}{growth} $\to$ \omterm{def:g11:dna-replication:helix}{DNA} copy $\to$ segregation protects ploidy.
\end{proposition}
\admitted

\section{Chromosome structure at mitosis}

\begin{definition}[Sister chromatids and centromere]\label{def:g11:cell-cycle-mitosis:chromatid}
After \omterm{def:g11:cell-cycle-mitosis:interphase}{S phase}, each \omterm{def:g11:dna-replication:chromosome}{chromosome} consists of two identical
\emph{sister chromatids}\index{sister chromatid} joined at a constricted region,
the \emph{centromere}\index{centromere}. \omterm{def:g10:molecules-of-life:proteins}{Protein} complexes called
\emph{kinetochores}\index{kinetochore} assemble at the centromere and attach to
spindle microtubules.
\end{definition}

\begin{definition}[Diploid chromosome number]\label{def:g11:cell-cycle-mitosis:diploid}
A \emph{diploid}\index{diploid} \omterm{def:g10:the-cell:cell}{cell} has two homologous sets of \omterm{def:g11:dna-replication:chromosome}{chromosomes}
($2n$). Humans: $2n = 46$. \omterm{def:g11:cell-cycle-mitosis:mitosis}{Mitosis} produces two daughter nuclei that are
genetically equivalent to the parent nucleus (apart from rare mutation), each
still $2n$ after chromatids separate and become \omterm{def:g11:dna-replication:chromosome}{chromosomes} of the daughters.
\end{definition}

\begin{omfigure}
\begin{tikzpicture}[font=\small]
  % unreplicated
  \draw[omDef, line width=3pt] (0,0) -- (0,2.2);
  \fill[omProp] (0,1.1) circle (0.12);
  \node[below, align=center, font=\footnotesize] at (0,-0.15)
    {before S\\(one DNA molecule)};
  % replicated
  \draw[omDef, line width=3pt] (3.2,0) -- (3.2,2.2);
  \draw[omMeth, line width=3pt] (3.7,0) -- (3.7,2.2);
  \fill[omProp] (3.45,1.1) circle (0.14);
  \draw[omProp, thick] (3.2,1.1) -- (3.7,1.1);
  \node[below, align=center] at (3.45,-0.15) {after S\\(two sister chromatids)};
  \node[above] at (3.2,2.25) {chromatid};
  \node[above] at (3.7,2.25) {chromatid};
  \node[right] at (3.9,1.1) {centromere};
\end{tikzpicture}

{\small One \omterm{def:g11:dna-replication:chromosome}{chromosome} before and after \omterm{def:g11:dna-replication:semi}{DNA replication}. Count \omterm{def:g11:dna-replication:chromosome}{chromosomes} by
\omterm{def:g11:cell-cycle-mitosis:chromatid}{centromeres}: still one \omterm{def:g11:dna-replication:chromosome}{chromosome} with two chromatids until anaphase separates
them.}
\end{omfigure}

\section{Mitosis: dividing the nucleus}

\begin{definition}[Mitosis]\label{def:g11:cell-cycle-mitosis:mitosis}
\emph{Mitosis}\index{mitosis} is nuclear division that segregates \omterm{def:g11:cell-cycle-mitosis:chromatid}{sister
chromatids} into two nuclei. Classic light-microscope stages are prophase,
prometaphase, metaphase, anaphase and telophase (names vary slightly by text;
the events matter more than the labels).
\end{definition}

\begin{proposition}[Events of mitosis]\label{prop:g11:cell-cycle-mitosis:stages}
\begin{itemize}
  \item \emph{Prophase}\index{prophase}: \omterm{def:g11:dna-replication:chromosome}{chromosomes} condense; the mitotic
        spindle begins to form; in animals the centrosomes move apart.
  \item \emph{Prometaphase}\index{prometaphase}: the \omterm{def:g10:the-cell:nucleus}{nuclear envelope} breaks
        down; microtubules attach to \omterm{def:g11:cell-cycle-mitosis:chromatid}{kinetochores}.
  \item \emph{Metaphase}\index{metaphase}: \omterm{def:g11:dna-replication:chromosome}{chromosomes} align at the spindle
        equator (metaphase plate).
  \item \emph{Anaphase}\index{anaphase}: cohesin between sisters is cleaved;
        \omterm{def:g11:cell-cycle-mitosis:chromatid}{sister chromatids} separate and move to opposite poles (each is now a
        daughter \omterm{def:g11:dna-replication:chromosome}{chromosome}).
  \item \emph{Telophase}\index{telophase}: \omterm{def:g11:dna-replication:chromosome}{chromosomes} decondense; \omterm{def:g10:the-cell:nucleus}{nuclear
        envelopes} re-form around the two sets.
\end{itemize}
\end{proposition}
\admitted

\begin{definition}[Cytokinesis]\label{def:g11:cell-cycle-mitosis:cytokinesis}
\emph{Cytokinesis}\index{cytokinesis} is division of the cytoplasm. Animal \omterm{def:g10:the-cell:cell}{cells}
pinch with a contractile ring of actin and myosin; plant \omterm{def:g10:the-cell:cell}{cells} build a new \omterm{prop:g11:cell-cycle-mitosis:cyto-compare}{cell
plate} (future \omterm{def:g10:the-cell:support}{cell wall}) between daughter nuclei. \omterm{def:g11:cell-cycle-mitosis:mitosis}{Mitosis} without cytokinesis
yields a multinucleate \omterm{def:g10:the-cell:cell}{cell}.
\end{definition}

\begin{proposition}[Animal vs plant cytokinesis]
\label{prop:g11:cell-cycle-mitosis:cyto-compare}
\begin{itemize}
  \item \textbf{Animal \omterm{def:g10:the-cell:cell}{cells}:} a contractile ring of actin and myosin tightens
        at the equator, forming a \emph{cleavage furrow}\index{cleavage furrow}
        that pinches the \omterm{def:g10:the-cell:cell}{cell} in two from the outside in.
  \item \textbf{Plant \omterm{def:g10:the-cell:cell}{cells}:} a rigid \omterm{def:g10:the-cell:support}{cell wall} prevents pinching. Golgi-derived
        vesicles fuse at the midplane to build a \emph{cell plate}\index{cell
        plate} that grows outward and matures into a new middle lamella and
        \omterm{def:g10:the-cell:support}{cell wall} separating the daughters.
\end{itemize}
Nuclear \omterm{def:g11:cell-cycle-mitosis:mitosis}{mitosis} is broadly similar; \omterm{def:g11:cell-cycle-mitosis:cytokinesis}{cytokinesis} machinery matches \omterm{def:g10:the-cell:support}{cell wall}
presence.
\end{proposition}
\admitted

\begin{omfigure}
\begin{tikzpicture}[font=\small]
  % metaphase cell
  \draw[omExo, thick] (0,0) ellipse (1.6 and 1.1);
  \foreach \x in {-0.6,-0.2,0.2,0.6}{
    \draw[omDef, line width=1.8pt] (\x,-0.35) -- (\x,0.35);
    \draw[omMeth, line width=1.8pt] (\x+0.12,-0.35) -- (\x+0.12,0.35);
  }
  \draw[omProp, dashed] (-1.3,0) -- (1.3,0);
  \node[below] at (0,-1.35) {metaphase};
  % anaphase
  \draw[omExo, thick] (4.2,0) ellipse (1.6 and 1.1);
  \foreach \x in {-0.5,-0.1,0.3,0.7}{
    \draw[omDef, line width=1.8pt] (4.2+\x,0.35) -- (4.2+\x,0.75);
    \draw[omMeth, line width=1.8pt] (4.2+\x,-0.35) -- (4.2+\x,-0.75);
  }
  \draw[-Stealth, omThm] (4.2,0.2) -- (4.2,0.85);
  \draw[-Stealth, omThm] (4.2,-0.2) -- (4.2,-0.85);
  \node[below] at (4.2,-1.35) {anaphase};
\end{tikzpicture}

{\small \omterm{prop:g11:cell-cycle-mitosis:stages}{Metaphase}: sisters aligned and bi-oriented. \omterm{prop:g11:cell-cycle-mitosis:stages}{Anaphase}: sisters separate
toward opposite poles --- equal segregation of replicated \omterm{prop:g11:dna-replication:genome}{genomes}.}
\end{omfigure}

\begin{omfigure}
\begin{tikzpicture}[font=\footnotesize]
  % animal furrow
  \draw[omExo, thick] (0,0.3) .. controls (0.8,-0.5) and (1.6,-0.5) .. (2.4,0.3);
  \draw[omExo, thick] (0,0.3) .. controls (0.8,1.1) and (1.6,1.1) .. (2.4,0.3);
  \draw[omThm, thick] (1.2,0.85) -- (1.2,-0.25);
  \node[align=center] at (1.2,-1.0) {animal:\\cleavage furrow};
  % plant cell plate
  \draw[omMeth, thick] (4.2,0) rectangle (6.8,1.4);
  \draw[omProp, very thick] (5.5,0.15) -- (5.5,1.25);
  \foreach \y in {0.35,0.7,1.05}
    \fill[omDef] (5.5,\y) circle [radius=0.06cm];
  \node[align=center] at (5.5,-0.55) {plant:\\cell plate from vesicles};
\end{tikzpicture}

{\small \omterm{def:g11:cell-cycle-mitosis:cytokinesis}{Cytokinesis} compared: outside-in pinching in animals versus inside-out
cell-plate construction in plants.}
\end{omfigure}

\begin{method}[Chromosome accounting through mitosis]\label{met:g11:cell-cycle-mitosis:count}
For a \omterm{def:g11:cell-cycle-mitosis:diploid}{diploid} \omterm{def:g10:the-cell:cell}{cell} with $2n = 4$ (two pairs):
\begin{enumerate}
  \item G$_1$: $4$ \omterm{def:g11:dna-replication:chromosome}{chromosomes}, $4$ \omterm{def:g11:dna-replication:helix}{DNA} molecules.
  \item G$_2$: $4$ \omterm{def:g11:dna-replication:chromosome}{chromosomes}, $8$ \omterm{def:g11:dna-replication:helix}{DNA} molecules (each \omterm{def:g11:dna-replication:chromosome}{chromosome} = $2$ chromatids).
  \item \omterm{prop:g11:cell-cycle-mitosis:stages}{Anaphase}: $8$ \omterm{def:g11:dna-replication:chromosome}{chromosomes} moving (sisters now separate).
  \item Each daughter in G$_1$: $4$ \omterm{def:g11:dna-replication:chromosome}{chromosomes}, $4$ \omterm{def:g11:dna-replication:helix}{DNA} molecules again.
\end{enumerate}
Always state whether you count \omterm{def:g11:dna-replication:chromosome}{chromosomes} (\omterm{def:g11:cell-cycle-mitosis:chromatid}{centromeres}) or chromatids/\omterm{def:g11:dna-replication:helix}{DNA}
molecules.
\end{method}

\section{Checkpoints: quality control}

\begin{definition}[Checkpoint]\label{def:g11:cell-cycle-mitosis:checkpoint}
A \emph{checkpoint}\index{checkpoint} is a control mechanism that can halt the
cycle until conditions are met (enough \omterm{def:g10:scientific-biology:properties}{growth}, intact \omterm{def:g11:dna-replication:helix}{DNA}, \omterm{def:g11:dna-replication:chromosome}{chromosomes} attached
to the spindle). Major checkpoints operate near the G$_1$/\omterm{def:g11:cell-cycle-mitosis:interphase}{S} transition, at the
G$_2$/M transition, and at the metaphase--anaphase transition (spindle assembly
checkpoint).
\end{definition}

\begin{proposition}[Role of checkpoints]\label{prop:g11:cell-cycle-mitosis:cp-role}
\omterm{def:g11:cell-cycle-mitosis:checkpoint}{Checkpoints} reduce the chance of propagating damaged \omterm{def:g11:dna-replication:helix}{DNA} or mis-segregated
\omterm{def:g11:dna-replication:chromosome}{chromosomes}. Failure of \omterm{def:g11:cell-cycle-mitosis:checkpoint}{checkpoint} control is a hallmark of many cancers: \omterm{def:g10:the-cell:cell}{cells}
divide despite \omterm{def:g11:dna-replication:helix}{DNA} damage or without proper attachment, generating genomic
instability.
\end{proposition}
\admitted

\begin{remark}
The molecular switches are cyclin-dependent kinases and their cyclin partners
--- names for later study. At high-school level, remember \emph{what} is checked
(\omterm{def:g10:scientific-biology:properties}{growth} signals, \omterm{def:g11:dna-replication:helix}{DNA} integrity, spindle attachment), not every \omterm{def:g10:molecules-of-life:proteins}{protein} name.
\end{remark}

\begin{example}[Spindle checkpoint logic]\label{ex:g11:cell-cycle-mitosis:spindle}
If one \omterm{def:g11:dna-replication:chromosome}{chromosome} is unattached at metaphase, the spindle \omterm{def:g11:cell-cycle-mitosis:checkpoint}{checkpoint} delays
anaphase. That pause allows time for microtubules to capture the free
\omterm{def:g11:cell-cycle-mitosis:chromatid}{kinetochore}, protecting against aneuploid daughters
(\cref{ch:g11:human-genetics}).
\end{example}

\begin{remark}[Cancer as checkpoint failure]
\label{rem:g11:cell-cycle-mitosis:cancer}
\emph{Cancer}\index{cancer} is uncontrolled \omterm{def:g10:the-cell:cell}{cell} proliferation. Many tumours
combine excessive \omterm{def:g10:scientific-biology:properties}{growth} signals with weakened \omterm{def:g11:cell-cycle-mitosis:checkpoint}{checkpoints}: \omterm{def:g10:the-cell:cell}{cells} may enter \omterm{def:g11:cell-cycle-mitosis:interphase}{S
phase} despite \omterm{def:g11:dna-replication:helix}{DNA} damage, or proceed through \omterm{def:g11:cell-cycle-mitosis:mitosis}{mitosis} with mis-attached
\omterm{def:g11:dna-replication:chromosome}{chromosomes}, generating genomic instability. Not every \omterm{def:g11:cell-cycle-mitosis:checkpoint}{checkpoint} failure is
cancer, and cancer involves more than the cycle alone (invasion, \omterm{def:g10:scientific-biology:properties}{metabolism},
immune evasion) --- but loss of cycle control is a central theme linking this
chapter to disease \omterm{def:g10:scientific-biology:science}{biology}.
\end{remark}

\begin{proposition}[Checkpoints protect the genome]
\label{prop:g11:cell-cycle-mitosis:genome-protect}
G$_1$/\omterm{def:g11:cell-cycle-mitosis:interphase}{S} \omterm{def:g11:cell-cycle-mitosis:checkpoint}{checkpoints} limit \omterm{def:g11:dna-replication:semi}{replication} of damaged \omterm{def:g11:dna-replication:helix}{DNA}; G$_2$/M \omterm{def:g11:cell-cycle-mitosis:checkpoint}{checkpoints} ask
whether \omterm{def:g11:dna-replication:semi}{replication} is complete; the spindle \omterm{def:g11:cell-cycle-mitosis:checkpoint}{checkpoint} asks whether sisters
are properly attached. Together they reduce the chance that daughter \omterm{def:g10:the-cell:cell}{cells}
inherit broken or mis-segregated \omterm{def:g11:dna-replication:chromosome}{chromosomes}.
\end{proposition}
\admitted

\section{Mitosis in the life of the organism}

\begin{proposition}[Where mitosis is used]\label{prop:g11:cell-cycle-mitosis:uses}
\omterm{def:g11:cell-cycle-mitosis:mitosis}{Mitosis} underlies \omterm{def:g10:scientific-biology:properties}{growth}, tissue renewal (skin, blood, gut lining), and asexual
\omterm{def:g10:scientific-biology:properties}{reproduction} in many \omterm{def:g10:scientific-biology:organism}{organisms}. It is not the division that halves \omterm{def:g11:dna-replication:chromosome}{chromosome}
number for sexual \omterm{def:g10:scientific-biology:properties}{reproduction} --- that is meiosis
(\cref{ch:g11:meiosis-life-cycles}).
\end{proposition}
\admitted

\section{Exercises}

\begin{exercise}[$\star$]\label{exo:g11:cell-cycle-mitosis:1}
Name the three stages of \omterm{def:g11:cell-cycle-mitosis:cycle}{interphase} and state what happens to \omterm{def:g11:dna-replication:helix}{DNA} in \omterm{def:g11:cell-cycle-mitosis:interphase}{S phase}.
\end{exercise}

\begin{exercise}[$\star$]\label{exo:g11:cell-cycle-mitosis:2}
What is the difference between a \omterm{def:g11:cell-cycle-mitosis:chromatid}{sister chromatid} and a homologous \omterm{def:g11:dna-replication:chromosome}{chromosome}?
\end{exercise}

\begin{exercise}[$\star$]\label{exo:g11:cell-cycle-mitosis:3}
List the stages of \omterm{def:g11:cell-cycle-mitosis:mitosis}{mitosis} in \omterm{def:g10:biodiversity-classification:hierarchy}{order} and give one defining event of metaphase
and of anaphase.
\end{exercise}

\begin{exercise}[$\star\star$]\label{exo:g11:cell-cycle-mitosis:4}
A human \omterm{def:g10:the-cell:cell}{cell} in G$_1$ has $46$ \omterm{def:g11:dna-replication:chromosome}{chromosomes}. How many \omterm{def:g11:dna-replication:chromosome}{chromosomes} and how many
chromatids does it have in G$_2$? In each daughter \omterm{def:g10:the-cell:cell}{cell} after \omterm{def:g11:cell-cycle-mitosis:mitosis}{mitosis} and
\omterm{def:g11:cell-cycle-mitosis:cytokinesis}{cytokinesis}?
\end{exercise}

\begin{exercise}[$\star\star$]\label{exo:g11:cell-cycle-mitosis:5}
Distinguish \omterm{def:g11:cell-cycle-mitosis:mitosis}{mitosis} from \omterm{def:g11:cell-cycle-mitosis:cytokinesis}{cytokinesis}. Can one occur without the other?
\end{exercise}

\begin{exercise}[$\star\star$]\label{exo:g11:cell-cycle-mitosis:6}
Why is the metaphase plate important for equal segregation of genetic material?
\end{exercise}

\begin{exercise}[$\star\star$]\label{exo:g11:cell-cycle-mitosis:7}
Name three \omterm{def:g11:cell-cycle-mitosis:checkpoint}{checkpoints} and one condition each is thought to monitor.
\end{exercise}

\begin{exercise}[$\star\star$]\label{exo:g11:cell-cycle-mitosis:8}
Plant \omterm{def:g10:the-cell:cell}{cells} lack a contractile ring like animal \omterm{def:g10:the-cell:cell}{cells}. How do they complete
\omterm{def:g11:cell-cycle-mitosis:cytokinesis}{cytokinesis}?
\end{exercise}

\begin{exercise}[$\star\star\star$]\label{exo:g11:cell-cycle-mitosis:9}
A drug prevents microtubule formation. Which mitotic events fail, and why might
the spindle \omterm{def:g11:cell-cycle-mitosis:checkpoint}{checkpoint} arrest the \omterm{def:g10:the-cell:cell}{cell}?
\end{exercise}

\begin{exercise}[$\star\star\star$]\label{exo:g11:cell-cycle-mitosis:10}
For $2n = 6$, sketch (words or numbers) \omterm{def:g11:dna-replication:chromosome}{chromosome} and chromatid counts in
G$_1$, G$_2$, anaphase, and G$_1$ of a daughter \omterm{def:g10:the-cell:cell}{cell}.
\end{exercise}

\begin{exercise}[$\star\star$]\label{exo:g11:cell-cycle-mitosis:11}
Compare \omterm{def:g11:cell-cycle-mitosis:cytokinesis}{cytokinesis} in a typical animal \omterm{def:g10:the-cell:cell}{cell} and a typical plant \omterm{def:g10:the-cell:cell}{cell}. Why
can plant \omterm{def:g10:the-cell:cell}{cells} not rely on a \omterm{prop:g11:cell-cycle-mitosis:cyto-compare}{cleavage furrow} alone?
\end{exercise}

\begin{exercise}[$\star\star\star$]\label{exo:g11:cell-cycle-mitosis:12}
Explain how failure of the G$_1$/\omterm{def:g11:cell-cycle-mitosis:interphase}{S} DNA-damage \omterm{def:g11:cell-cycle-mitosis:checkpoint}{checkpoint} and failure of the
spindle \omterm{def:g11:cell-cycle-mitosis:checkpoint}{checkpoint} could each promote tumour \omterm{def:g10:scientific-biology:properties}{development}. Why is ``faster
\omterm{def:g11:cell-cycle-mitosis:mitosis}{mitosis}'' alone an incomplete description of cancer?
\end{exercise}

\begin{problem}[Following one chromosome through a cycle]\label{pb:g11:cell-cycle-mitosis:1}
Consider a \omterm{def:g11:cell-cycle-mitosis:diploid}{diploid} animal \omterm{def:g10:the-cell:cell}{cell} with $2n = 4$ \omterm{def:g11:dna-replication:chromosome}{chromosomes} (two homologous pairs).
\begin{enumerate}
  \item Draw the \omterm{def:g10:the-cell:cell}{cell} in G$_1$ with two long and two short \omterm{def:g11:dna-replication:chromosome}{chromosomes} (homologue
        pairs). Mark maternal and paternal homologues differently.
  \item Show the same \omterm{def:g10:the-cell:cell}{cell} in G$_2$ after \omterm{def:g11:dna-replication:semi}{replication}.
  \item Describe \omterm{def:g11:dna-replication:chromosome}{chromosome} behaviour from prophase through telophase for one
        \omterm{def:g11:dna-replication:chromosome}{chromosome} (condensation, alignment, separation).
  \item How many \omterm{def:g11:dna-replication:chromosome}{chromosomes} does each daughter \omterm{def:g10:the-cell:cell}{cell} receive? Are the daughters
        genetically identical to the parent (ignoring mutation)?
  \item If the spindle \omterm{def:g11:cell-cycle-mitosis:checkpoint}{checkpoint} fails and both sisters of one \omterm{def:g11:dna-replication:chromosome}{chromosome} go to
        the same pole, what are the \omterm{def:g11:dna-replication:chromosome}{chromosome} numbers of the two nuclei?
  \item Why must \omterm{def:g11:dna-replication:semi}{DNA replication} be completed before anaphase begins?
  \item Name one tissue in a mammal where \omterm{def:g11:cell-cycle-mitosis:mitosis}{mitosis} is frequent in adults, and one
        where most \omterm{def:g10:the-cell:cell}{cells} remain in G$_0$.
\end{enumerate}
\end{problem}
